\section{A downstairs search on the tetraquadric quotient}
\label{sec:downstairs}

\subsection{Wilson lines and equivariant structures}

For abelian $\Gam$ the Wilson line is a pair of characters: $\chi_T$ on the colour
triplet and $\chi_D$ on the weak doublet of the fundamental $\five$. Everything
else follows from $\ten=\wedge^2\five$:

\begin{center}\small
\begin{tabular}{@{}llll@{}}
\toprule
multiplet & pieces & Wilson-line character & $6Y$ \\
\midrule
$\ten$ & $Q$, $u^c$, $e^c$ & $\chi_T\chi_D$, $\chi_T^2$, $\chi_D^2$ & $1,-4,6$ \\
$\bfive$ & $d^c$, $L$ & $\chi_T^{-1}$, $\chi_D^{-1}$ & $2,-3$ \\
\bottomrule
\end{tabular}
\end{center}

\begin{lemma}\label{lem:wilson}
The Wilson line lies in $SU(5)$, so $\chi_T^3\chi_D^2=1$. For a group of exponent
two this forces $\chi_T=1$ and leaves $\chi_D$ free.
\end{lemma}

A piece with Wilson-line character $\chi$ on a bundle with equivariant structure
$e$ survives with the multiplicity of the upstairs character $\chi-e$ (in
additive notation). Equivariant structures are enumerated subject to
$\sum_ae_a=0$, so that $\det V$ descends. The Wilson line is thus derived, not
chosen; for $\Xa/(\Z_2\times\Z_2)$ it reproduces the convention of the published
database.

\subsection{Results}

Starting from the 58 ambient-row orbits of our exhaustive $\Gam=\Z_2\times\Z_2$,
$K=5$ scan of $\Xa$, with no published input beyond the symmetry data:

\begin{center}\small
\begin{tabular}{@{}lr@{}}
\toprule
upstairs orbits examined & 58 \\
\quad some $L_a$ not equivariant ($\sum_ik^i$ odd, \cref{prop:equivariance}) & 26 \\
\quad equivariant, all needed cohomology determined & 32 \\
\qquad no Wilson line and equivariant structure gives an acceptable spectrum & 11 \\
\qquad \textbf{three chiral families, no charged chiral exotics, a Higgs pair} & \textbf{21} \\
\bottomrule
\end{tabular}
\end{center}

All 10 orbits of the published standard-model database on this quotient are
recovered. The other 11 are absent from it. We call them \emph{candidates} and
label them A01--A11.

\subsection{Audit of the eleven candidates}

Each candidate was frozen into an immutable manifest and audited against every
condition of the original construction:

\begin{center}\small
\begin{tabular}{@{}lr@{}}
\toprule
condition & result \\
\midrule
$c_1(V)=0$, structure group, anomaly, Bogomolov & 11/11 \\
poly-stability with an \emph{exact} witness & 11/11 (8 rational, 3 real-algebraic) \\
$\ind V=\ind\wedge^2V=-12$ & 11/11 \\
$h^\bullet(V)=(0,12,0,0)$, $h^\bullet(\wedge^2V)=(0,12+n_h,n_h,0)$, $n_h>0$ & 11/11 \\
every bundle equivariant; all character tables determined & 11/11 \\
three $Q,u^c,d^c,L,e^c$; no charged chiral exotics; a Higgs pair & 11/11 \\
no colour-triplet Higgs & 11/11 \\
absent from every published downstairs model & 11/11 \\
\bottomrule
\end{tabular}
\end{center}

For charge matrices of rank two the slope-zero locus is a curve, so rationalising
a numerical point cannot produce an exact witness; a rational point \emph{on}
the curve is found instead (for A02, $t=(15,15,3,1)$). Two of the real-algebraic
witnesses have $t_1=4+\sqrt{21}$.

\subsection{Why the candidates are absent from the published database}

The published standard-model database was built from a GUT scan over
$-3\le k_a^i\le3$ on $h^{1,1}=4$ \cite[Sec.~5.3]{AGLP2012}, which additionally
demanded $h^1(X,L_a^*\otimes L_b^*)<|\Gam|$ for every pair ``to increase the
chance that the Higgs triplets can be removed''~\cite{AGLP2011}. We checked both
conditions against the data: every published model satisfies them (2,111 of
2,111 with readable cohomology for the second). Every candidate violates at
least one:

\begin{center}\footnotesize
\setlength{\tabcolsep}{4pt}
\begin{tabular}{@{}lcccccl@{}}
\toprule
 & $\max|k|$ & extra $U(1)$s & $n_h$ & $\max h^1(L_a^*L_b^*)$ & configurations & outside the scan by \\
\midrule
A01 & 5 & 1 & 3 & 3 & 192 & charge range \\
A02 & 5 & 2 & 17 & 8 & 336 & charge range, $h^1$ bound \\
A03 & 5 & 1 & 3 & 3 & 192 & charge range \\
A04 & 4 & 2 & 17 & 8 & 336 & charge range, $h^1$ bound \\
A05 & 4 & 1 & 3 & 3 & 192 & charge range \\
A06 & 4 & 2 & 17 & 8 & 336 & charge range, $h^1$ bound \\
A07 & 3 & 2 & 17 & 8 & 336 & $h^1$ bound \\
A08 & 3 & 1 & 11 & 8 & 192 & $h^1$ bound \\
A09 & 3 & 1 & 11 & 8 & 192 & $h^1$ bound \\
A10 & 3 & 2 & 17 & 8 & 336 & $h^1$ bound \\
A11 & 3 & 2 & 17 & 8 & 336 & $h^1$ bound \\
\bottomrule
\end{tabular}
\end{center}

So the candidates lie outside the published scan's \emph{declared} domain; they
were not missed by it. Two remarks. The $h^1$ bound is a device to make
doublet--triplet splitting likely, not a physical requirement: our
computation shows that all eleven nevertheless admit 192--336 equivariant
structures and Wilson lines each that keep a Higgs doublet pair and project out
every colour triplet. And the six rank-two candidates carry two
additional massless $U(1)$s where every published model on this quotient carries
one, so they are phenomenologically distinct rather than mere additions to the
catalog. A01, A03 and A05 are profile-identical to published models---rank
three, one extra $U(1)$, $n_h=3$---and differ only in charge range.

We stress what the audit does not establish: that these are viable. That is what
the Yukawa analysis of the next section addresses.
